\documentclass[11pt,a4paper]{article}
\usepackage[margin=1in]{geometry}
\usepackage{amsmath,amssymb,amsfonts}
\usepackage{mathtools}
\usepackage{lmodern}
\usepackage[T1]{fontenc}
\usepackage[utf8]{inputenc}
\usepackage{textcomp}
\usepackage{microtype}
\usepackage{parskip}
\usepackage{enumitem}
\usepackage{array}
\usepackage{booktabs}
\usepackage{hyperref}
\hypersetup{hidelinks,
  pdftitle={Paper 095 -- Form-FORCED / Value-INHERITED Methodology},
  pdfauthor={Allen Proxmire}}
\setlength{\parindent}{0pt}
\setlength{\parskip}{6pt plus 2pt minus 1pt}

\title{\vspace{-1cm}\Large\bfseries Paper 095 --- Form-FORCED / Value-INHERITED Methodology}
\author{Allen Proxmire}
\date{May 2026}

\begin{document}
\maketitle

\section*{Preamble --- What This Paper Does NOT Claim}

\begin{enumerate}[noitemsep]
\item This paper does \textbf{not} claim novel substantive content beyond formalizing the corpus's methodological discipline.
\item It does \textbf{not} override any constituent-paper verdict.
\item It does \textbf{not} introduce new substrate primitives.
\item ``Form-FORCED / value-INHERITED methodology'' = the corpus's consistent rule: structural form of results is FORCED by substrate primitives + declared postulates; numerical values, coefficients, and matching parameters are INHERITED from empirical / standard-physics literature.
\end{enumerate}

\section*{Abstract}

The ED corpus employs a consistent \textbf{form-FORCED / value-INHERITED} methodology: across QM-emergence (Phase-1), QFT-arc, gravity arc, BH arc, Q-Compute arc, entanglement arc, soft-matter arc, and wedges arc, structural forms (equations, inequalities, classifications, scaling laws) are FORCED by substrate primitives + declared postulates; numerical values, coefficients, and matching parameters are INHERITED. This paper formalizes the methodology, supplies a verdict-classification grammar (FORCED-unconditional / Intermediate Path C / form-FORCED + value-INHERITED), and identifies the methodology's structural strengths and limits.

\section{Primitive Inputs and Paper-Specific Postulates}

\subsection{Primitives invoked (per Paper\_087)}

This is a methodology paper; primitives are invoked only at the meta-level. The 13 primitives P01--P13 are the substrate-level content; methodology is what the corpus does with them.

\subsection{Upstream dependencies}

\begin{itemize}[noitemsep]
\item \textbf{I-087:} Paper\_087 13-primitive position paper.
\item \textbf{I-088:} Paper\_088 primitive load-bearing audit.
\item \textbf{I-089:} Paper\_089 V1 canonical reference (exemplar of methodology applied).
\end{itemize}

\subsection{Paper-specific postulates}

\begin{itemize}[noitemsep]
\item \textbf{P-Methodology-Triple:} The corpus methodology classifies all results into one of three verdict-tiers:
\begin{itemize}[noitemsep]
\item \textbf{(M1) FORCED-unconditional:} result follows from the 13 primitives + standard math, no additional paper-specific postulates required.
\item \textbf{(M2) Intermediate Path C:} result follows from primitives + standard math + a small number of declared paper-specific postulates whose constructive derivation is OPEN.
\item \textbf{(M3) form-FORCED + value-INHERITED:} structural form follows under primitives + postulates; numerical values inherited via empirical matching.
\end{itemize}
\item \textbf{P-Methodology-Consistency:} All Wave-2 papers across the corpus apply this methodology consistently.
\end{itemize}

\section*{\S2.5 Audit Table}

\begin{center}
\begin{tabular}{cllp{0.45\textwidth}}
\toprule
\textbf{\#} & \textbf{Step} & \textbf{Label} & \textbf{Notes} \\
\midrule
1 & 13 primitives & I (axiomatic) & Paper\_087. \\
2 & Primitive load-bearing audit & I & Paper\_088. \\
3 & Three-tier verdict classification & P-Methodology-Triple & Postulate. \\
4 & Cross-corpus consistent application & P-Methodology-Consistency & Postulate. \\
5 & Specific verdict assignments & I & Per-paper inheritance. \\
\bottomrule
\end{tabular}
\end{center}

\section{The Methodology}

\subsection{Three verdict tiers}

\begin{itemize}[noitemsep]
\item \textbf{(M1) FORCED-unconditional.} Example: Phase-1 QM emergence (four QM postulates forced from primitives + standard math; no additional postulates declared at the QM-postulate level).
\item \textbf{(M2) Intermediate Path C.} Example: NS-smoothness (Paper\_077), YM Clay-relevance (Paper\_023), Arc ED-10 (Paper\_033), NS-MHD H1/H2/H3 closure (Paper\_081). Substrate mechanisms identified; postulates declared; constructive derivations OPEN.
\item \textbf{(M3) form-FORCED + value-INHERITED.} Example: Newton's $G$ (Paper\_027; form FORCED, $G$ value INHERITED), BTFR slope-4 (Paper\_031), area-law (Paper\_043), Krieger--Dougherty rheology (Paper\_079).
\end{itemize}

\subsection{Honest labels}

The methodology requires every load-bearing step in every paper to be labeled:
\begin{itemize}[noitemsep]
\item \textbf{P} (postulate / definition).
\item \textbf{A} (analogy / regime / position).
\item \textbf{I} (inherited / empirical / standard math).
\item \textbf{D} (actual derivation).
\item \textbf{D-via-I} (application of standard machinery to substrate-adapted content under declared postulates).
\item \textbf{OPEN} (not claimed; flagged for future work).
\end{itemize}

\subsection{No conflation}

Standard-form constructions (e.g., writing down the Klein-Gordon equation, the Lindblad form, the MOND-like field equation, the tensor product $\Psi^A \otimes \Psi^B$) are \textbf{P (definitional)}, not D. Empirically matched parameters are \textbf{I}, not D. Composite verdicts assembled from upstream inheritances are \textbf{D} (the composition itself is new content) or \textbf{A$\to$position} (when the verdict is a framing claim).

\subsection{Methodological strengths}

\begin{itemize}[noitemsep]
\item \textbf{Falsifiability:} every postulate has a stated F-line.
\item \textbf{Honest cumulativity:} future derivations strengthen by reducing P-postulates to D-rows.
\item \textbf{No overclaim:} numerical values are not claimed as derived unless they are.
\end{itemize}

\subsection{Methodological limits}

\begin{itemize}[noitemsep]
\item Many corpus results are at \textbf{(M3)} level: form derived, values not. This is a strength relative to ``derive everything'' overclaims but a limit relative to constructive-QFT or constructive-relativity ambitions.
\item The methodology does not produce constructive proofs of Clay-class problems; it produces Intermediate Path C verdicts (M2).
\end{itemize}

\section{Falsification Criteria}

\begin{itemize}[noitemsep]
\item \textbf{F1:} Discovery of a corpus paper where load-bearing steps are unlabeled or mislabeled at scale --- refutes P-Methodology-Consistency.
\item \textbf{F2:} Identification of a verdict that does not fit any of the three tiers (M1, M2, M3) --- refutes P-Methodology-Triple's exhaustiveness.
\item \textbf{F3:} Discovery of a corpus result claimed as D but actually I, P, or A on rereading --- falsifies the specific paper's discipline.
\end{itemize}

\section{Position Statement}

The ED corpus employs a \textbf{form-FORCED / value-INHERITED methodology} with three-tier verdict classification (M1 / M2 / M3) under P-Methodology-Triple + P-Methodology-Consistency. The methodology is \textbf{the academic case}: cross-domain reach + honest labels.

\end{document}
